\documentclass[10pt, aspectratio=169]{beamer}

\usepackage[spanish]{babel}
\usepackage[utf8]{inputenc}
\usepackage[pscoord]{eso-pic}
\usepackage{hyperref}
\usepackage{eqnarray,amsmath}
\usepackage[absolute,overlay]{textpos}
\hypersetup{colorlinks=true,urlcolor=blue}

\usepackage{tikz}
\usetikzlibrary{automata, positioning, arrows.meta}

\usepackage{listings}
\usepackage{xcolor}

\usepackage[style=apa]{biblatex}
\bibliography{biblio.bib}

\title{}
\author{}
\date{}

\newtheorem{teorema}{Teorema.}
\setbeamersize{text margin left=10pt,text margin right=10pt}
\beamertemplatenavigationsymbolsempty

\hypersetup{
	colorlinks=true,
	pdfborderstyle={/S/U/W 1}
}

\setbeamertemplate{background}{\includegraphics[width=\paperwidth,height=\paperheight,keepaspectratio]{portada.pdf}}

\newcommand{\azulULA}{\color[rgb]{0.30078125,0.30078125,0.999999}}
\definecolor{azulULA}{RGB}{77, 77, 255}

\setbeamercolor{title}{fg=azulULA}
\setbeamercolor{frametitle}{fg=azulULA}
\setbeamercolor{block title}{bg=azulULA,fg=white}
\setbeamercolor{block body}{bg=azulULA!10!bg}

\renewcommand{\qedsymbol}{\textcolor{black}{\openbox}}

\newcommand{\placetextbox}[3]{
	\setbox0=\hbox{#3}
	\AddToShipoutPictureFG*{
		\put(\LenToUnit{#1\paperwidth},\LenToUnit{#2\paperheight}){\vtop{{\null}\makebox[0pt][c]{#3}}}
	}
}
\newcommand{\placetextboxleft}[3]{
	\setbox0=\hbox{#3}
	\AddToShipoutPictureFG*{
		\put(\LenToUnit{#1\paperwidth},\LenToUnit{#2\paperheight}){\vtop{{\null}\makebox[0pt][l]{#3}}}
	}
}

\begin{document}
	% Diapositiva de la portada.
	\begin{frame}
		\placetextbox{0.47}{0.40}{\large\azulULA\textbf{\begin{tabular}{c} Teoría de la Computación y Lenguajes Formales: \\
					Sesión 10. Evolución del Motor a PDA y Seguimiento de Workflows \end{tabular}}}
	\end{frame}
	
	% Cambio de formato para las diapositivas centrales (cuerpo/contenido).
	\setbeamertemplate{background}{
		\includegraphics[width=\paperwidth,height=\paperheight,keepaspectratio]{fondo1280x720.pdf}}
	
	\begin{frame}{Alcance de la sesión 10}
		\tableofcontents
	\end{frame}

% =======================================================
\section{De DFA a PDA: ¿Por qué evoluciona el Motor?}
% =======================================================

\begin{frame}{}
	\centering
	\Huge{Evolución de la Arquitectura: \\ De DFA Plano a PDA con Memoria Estructurada}
\end{frame}

\begin{frame}{El Límite del Motor DFA de la Sesión 7}
	\begin{block}{¿Cómo operaba el motor DFA básico?}
		En el Hito 1 programamos un motor que evaluaba transiciones directas basadas únicamente en el estado actual y el evento de entrada:
		\[ \delta(\text{estado\_actual}, \, \text{evento}) \;\longrightarrow\; \text{nuevo\_estado} \]
	\end{block}
	
	\begin{alertblock}{La necesidad en Workflows Complejos}
		Cuando un flujo de trabajo incluye **sub-procesos anidados o fases jerárquicas**:
		\begin{itemize}
			\item El sistema debe **recordar el contexto del proceso padre** al ingresar a un sub-workflow.
			\item Debe garantizar que las fases abiertas se cierren en el **orden inverso estricto (LIFO)**.
			\item Debe validar que no queden tareas subordinadas pendientes antes de dar por completado el producto.
		\end{itemize}
	\end{alertblock}
\end{frame}

% =======================================================
\section{Extensión de la Ontología en YAML}
% =======================================================

\begin{frame}{}
	\centering
	\Huge{La Especificación del PDA: \\ Extensión de la Ontología en YAML}
\end{frame}

\begin{frame}[fragile]{Estructura del YAML para el PDA}
	\begin{block}{Declaración de la 7-Tupla en YAML}
		Añadimos a la especificación los elementos de la memoria de pila:
		\begin{lstlisting}[language=HTML, basicstyle=\tiny, keywordstyle=\color{blue}]
recurso: "Centro_Mecanizado_CNC"
tipo: "Estacion_Con_Subprocesos"
estado_inicial: "OCIOSO"
simbolo_inicial_pila: "Z0"
alfabeto_pila: ["Z0", "CTX_DESBASTE", "CTX_ACABADO"]

estados_finales:
  exito: ["COMPLETADO"]
  falla: ["BLOQUEO_CONTEXTO"]
		\end{lstlisting}
	\end{block}
	
	\begin{block}{Definición de Transiciones con Pila ($\delta: Q \times \Sigma \times \Gamma \to Q \times \Gamma^*$)}
		\begin{lstlisting}[language=HTML, basicstyle=\tiny, keywordstyle=\color{blue}]
transiciones_pda:
  # 1. PUSH: Iniciar sub-proceso de Desbaste (apila CTX_DESBASTE sobre Z0)
  - de: "OCIOSO"
    evento: "iniciar_desbaste"
    tope_pila: "Z0"
    a: "EN_DESBASTE"
    apilar: ["CTX_DESBASTE", "Z0"]

  # 2. POP: Finalizar sub-proceso de Desbaste (desapila CTX_DESBASTE)
  - de: "EN_DESBASTE"
    evento: "fin_desbaste"
    tope_pila: "CTX_DESBASTE"
    a: "DESBASTE_LISTO"
    apilar: []  # Cadena vacia (POP)
		\end{lstlisting}
	\end{block}
\end{frame}

% =======================================================
\section{Implementación del Motor PDA en Python}
% =======================================================

\begin{frame}{}
	\centering
	\Huge{Implementación del Motor en Python: \\ La Clase PDA y la Pila Física}
\end{frame}

\begin{frame}[fragile]{La Clase PDA: Estructura de Datos}
	\begin{block}{Inicialización con Pila LIFO}
		\begin{lstlisting}[language=Python, basicstyle=\tiny, keywordstyle=\color{blue}]
class PDA:
    def __init__(self, states, alphabet, stack_alphabet, transitions, q0, z0, final_states):
        self.states = states
        self.alphabet = alphabet
        self.stack_alphabet = stack_alphabet
        self.transitions = transitions  # Clave: (estado, evento, tope_pila) -> (nuevo_estado, push_list)
        self.q0 = q0
        self.z0 = z0
        self.final_states = final_states
        
        self.current_state = q0
        self.stack = [z0]  # Pila inicializada con el simbolo de fondo
        self.history = []

    def get_stack_top(self):
        return self.stack[-1] if self.stack else None
		\end{lstlisting}
	\end{block}
\end{frame}

\begin{frame}[fragile]{Lógica de Transición con Pila en Tiempo Real}
	\begin{block}{Método de Evolución $\delta(q, e, Z)$}
		\begin{lstlisting}[language=Python, basicstyle=\tiny, keywordstyle=\color{blue}]
    def transition(self, event):
        if event not in self.alphabet:
            raise ValueError(f"[Error Alfabeto] Evento '{event}' no reconocido en el catalogo.")
        
        stack_top = self.get_stack_top()
        if stack_top is None:
            raise RuntimeError("[Error Fatal] Pila vacia. Bloqueo de ejecucion.")

        key = (self.current_state, event, stack_top)
        
        if key not in self.transitions:
            raise ValueError(
                f"[Error de Transicion] No existe regla para ({self.current_state}, '{event}', tope={stack_top}). "
                f"El contexto activo no permite este evento."
            )

        new_state, push_symbols = self.transitions[key]
        
        # 1. POP del simbolo actual del tope
        popped = self.stack.pop()
        
        # 2. PUSH de los nuevos simbolos (en orden inverso para que push_symbols[0] quede arriba)
        for sym in reversed(push_symbols):
            self.stack.append(sym)
            
        prev_state = self.current_state
        self.current_state = new_state
        self.history.append((event, prev_state, new_state, self.stack.copy()))
        
        print(f"-> Evento '{event}': {prev_state} -> {new_state} | Pila: {self.stack}")
		\end{lstlisting}
	\end{block}
\end{frame}

% =======================================================
\section{Seguimiento a un Flujo de Trabajo (Tracking)}
% =======================================================

\begin{frame}{}
	\centering
	\Huge{Seguimiento de Workflows en Tiempo Real: \\ Detección de Anomalías y Contexto}
\end{frame}

\begin{frame}{La Arquitectura de Seguimiento: Sensores y Gemelo Digital}
	\begin{block}{El Sensor como Puente entre la Realidad y el Modelo}
		Para que el motor pueda realizar el seguimiento en tiempo real, los hechos del mundo físico u operacional deben convertirse en eventos discretos:
		\begin{itemize}
			\item \textbf{Sensor Automático (Físico / IoT):} Final de carrera, sensor fotoeléctrico, lector RFID/código de barras, señal de PLC.
			\item \textbf{Sensor Humano (Operador / Analista):} El operario que observa el estado de la pieza o el documento e introduce el evento en una terminal HMI/consola (\texttt{input()}).
		\end{itemize}
	\end{block}

	\begin{alertblock}{El Motor como Gemelo Digital (\emph{Digital Twin})}
		El motor PDA mantiene en memoria el \textbf{Gemelo Digital} del proceso. Cada vez que un sensor emite un evento $e \in \Sigma$:
		\[ \text{Mundo Real} \;\xrightarrow{\text{Hecho Físico}}\; \text{Sensor/HMI} \;\xrightarrow{\text{Evento } e}\; \text{Motor PDA (Gemelo en RAM)} \]
		El gemelo evoluciona su estado y su pila. Si el evento contradice las reglas del proceso o el contexto activo, el gemelo \textbf{detecta la discrepancia y emite una alarma}.
	\end{alertblock}
\end{frame}

\begin{frame}{Caso de Estudio: Mecanizado con Sub-procesos Anidados}
	\begin{block}{Jerarquía del Flujo de Trabajo}
		El producto debe pasar obligatoriamente por dos sub-procesos jerárquicos:
		\begin{enumerate}
			\item \textbf{Sub-proceso 1: Desbaste Grueso} (Apertura $\to$ PUSH \texttt{CTX\_DESBASTE}, Cierre $\to$ POP).
			\item \textbf{Sub-proceso 2: Acabado Fino} (Apertura $\to$ PUSH \texttt{CTX\_ACABADO}, Cierre $\to$ POP).
		\end{enumerate}
	\end{block}

	\begin{center}
		\begin{tikzpicture}[
			shorten >=1pt, 
			node distance=3.2cm, 
			on grid, 
			auto,
			>={Stealth[length=2.5mm]},
			every state/.style={draw=blue!70, fill=blue!5, thick, minimum size=1.2cm, font=\scriptsize}
			]
			\node[state, initial] (q0) {OCIOSO};
			\node[state] (q1) [right=of q0] {DESBASTE};
			\node[state] (q2) [right=of q1] {ACABADO};
			\node[state, accepting] (q3) [right=of q2] {LISTO};
			
			\path[->, thick] 
			(q0) edge node [above] {\tiny in\_desb, $Z_0 / D Z_0$} (q1)
			(q1) edge node [above] {\tiny fin\_desb, $D / \varepsilon$} (q2)
			(q2) edge node [above] {\tiny fin\_acab, $A / \varepsilon$} (q3);
		\end{tikzpicture}
	\end{center}
\end{frame}

\begin{frame}{Monitoreo y Detección de Errores por el Motor}
	\begin{block}{¿Cómo detecta el motor las anomalías operacionales?}
		\begin{itemize}
			\item \textbf{Caso 1: Ejecución Normal Balanceada}
			\[ \text{in\_desb} \;\to\; \text{fin\_desb} \;\to\; \text{in\_acab} \;\to\; \text{fin\_acab} \]
			La pila inicia en $[Z_0]$, crece en cada sub-fase, se vacía a $[Z_0]$ y llega al estado marcado $\Rightarrow$ \textbf{ACEPTADO}.
			\item \textbf{Caso 2: Intento de Salto de Fase (Violación de Contexto)}
			Si el operario intenta enviar \texttt{fin\_acab} estando en la fase de desbaste: el tope de la pila es \texttt{CTX\_DESBASTE}, por lo que la regla de acabado es rechazada de inmediato $\Rightarrow$ \textbf{ANOMALÍA DETECTADA}.
			\item \textbf{Caso 3: Tarea Inconclusa al Despachar}
			Si se envía el evento \texttt{despachar} mientras la pila contiene $[Z_0, \text{CTX\_ACABADO}]$, el motor rechaza la finalización porque la pila no regresó a su símbolo base $\Rightarrow$ \textbf{PRODUCTO BLOQUEADO}.
		\end{itemize}
	\end{block}
\end{frame}

\begin{frame}{Caso de Negocio: Proceso de Compras entre 4 Dependencias}
	\begin{block}{El Gemelo Digital en Procesos Administrativos (ERP / BPMN)}
		Un pedido de compra transita por 4 dependencias (Solicitante $\to$ Compras $\to$ Almacén $\to$ Finanzas). Los analistas actúan como ``sensores humanos'' que registran eventos en el sistema:
	\end{block}

	\begin{center}
		\begin{tikzpicture}[
			shorten >=1pt, 
			node distance=2.9cm, 
			on grid, 
			auto,
			>={Stealth[length=2.5mm]},
			every state/.style={draw=blue!70, fill=blue!5, thick, minimum size=1.1cm, font=\tiny}
			]
			\node[state, initial] (q0) {SOLICITUD};
			\node[state] (q1) [right=of q0] {EVALUAC.};
			\node[state] (q2) [right=of q1] {RECEPCION};
			\node[state, accepting] (q3) [right=of q2] {PAGO};
			\node[state, accepting] (q_abort) [below=1.4cm of q1, draw=red!70, fill=red!5] {RECHAZO};
			
			\path[->, thick] 
			(q0) edge node [above] {\tiny solicitar} (q1)
			(q1) edge node [above] {\tiny orden\_compra} (q2)
			(q2) edge node [above] {\tiny recep\_conforme} (q3)
			(q1) edge node [left] {\tiny rechazo\_pedido} (q_abort)
			(q2) edge node [right] {\tiny rechazo\_producto} (q_abort);
		\end{tikzpicture}
	\end{center}

	\begin{alertblock}{Bifurcaciones que Cortan el Flujo}
		Los eventos de rechazo (\texttt{rechazo\_pedido} o \texttt{rechazo\_producto}) desvían inmediatamente el token hacia el estado de falla (\texttt{RECHAZO}), cortando el flujo e impidiendo que se genere la orden de pago.
	\end{alertblock}
\end{frame}

% =======================================================
\section{Conclusión}
% =======================================================

\begin{frame}{Conclusión de la Sesión 10}
	\begin{block}{Logros del Motor PDA}
		\begin{itemize}
			\item \textbf{Control de Contexto Dinámico:} La Pila permite al motor saber exactamente en qué nivel de sub-proceso se encuentra el producto.
			\item \textbf{Validación de Balances:} Garantiza que ninguna fase de trabajo quede abierta o se cierre fuera de orden.
			\item \textbf{Trazabilidad Total:} El historial de ejecuciones registra no solo los estados, sino la evolución de la memoria de contexto en cada evento.
		\end{itemize}
	\end{block}

	\begin{exampleblock}{Próximo Bloque (Semanas 10-12): Concurrencia y Redes de Petri}
		Ya dominamos el control de un recurso con sub-procesos anidados. En la próxima clase daremos el salto a la **concurrencia**: múltiples recursos comunicándose mediante **paso de mensajes ($!a / ?a$) y Redes de Petri**.
	\end{exampleblock}
\end{frame}

% Cambio de formato para la diapositiva final.
\setbeamertemplate{headline}{}
\setbeamertemplate{background}
{\includegraphics[width=\paperwidth,height=\paperheight,keepaspectratio]{salida.pdf}}

% Diapositiva final.
\begin{frame}
	\placetextbox{0.6}{0.87}{\Huge\textbf{¡Muchas gracias!}}
\end{frame}

\setbeamertemplate{headline}{}
\setbeamertemplate{background}{}

\printbibliography
\end{document}
